\section{Calibration, initial-state and residual assumptions}
\label{sec:nuisance-assumptions}

\subsection{Heater calibration introduces another exact ambiguity}
The conversion $c=\alpha P$ is fixed in the baseline study. If it is unknown,
the transformation
\[
 (c,\Ua,\Cp)\longmapsto(\lambda c,\lambda\Ua,\lambda\Cp),
 \qquad \lambda>0,
\]
preserves $K=c/\Ua$, $\tau=\Cp/\Ua$, and every reduced-model prediction under
the same initial conditions. In the original model, scaling $c,\Ua,\Ub,\CpH$
and $\CpS$ together has the same effect. Input design cannot identify this
absolute thermal scale from temperature and percentage command alone.

Let $C_{\mathrm{io}}$ be a covariance matrix for
$(\log K,\log\tau,\log\theta)$ and suppose an independent heater calibration
has log-variance $v_c$. With the transformation matrix $M$ defined in
Appendix~\ref{sec:structural}, the corresponding physical log-covariance is
\[
 C_{\mathrm{phys}}=M C_{\mathrm{io}}M^\top+v_c\bm v\bm v^\top,
 \qquad \bm v=(1,1,0)^\top.
\]
This covariance identity is exact in logarithmic coordinates under the stated
independence assumption; Gaussianity is not needed for the covariance itself.
It adds correlated uncertainty to conductance and capacity while leaving delay
unchanged. A hypothetical lognormal calibration with coefficient of variation
$r_c$ has $v_c=\log(1+r_c^2)$. Such a sensitivity calculation must be labelled
as assumed calibration uncertainty: these data do not supply $r_c$.

\subsection{Ambient and history are nuisance quantities}
For either prescribed delay-state initialization, hold the observed initial
sensor temperature $T_0$ fixed and vary ambient temperature. Linearity gives
\[
 \frac{\partial T(t)}{\partial\Tamb}=1-e^{-t/\tau}.
\]
An ambient error therefore persists rather than averaging away with additional
samples in one record. Under constant input $u_0$ and equilibrated delay states,
\[
 T(t)=T_0 e^{-t/\tau}+(\Tamb+Ku_0)(1-e^{-t/\tau}).
\]
Only the combination $\Tamb+Ku_0$ is identified from this response if ambient
is free, and delay is absent. Independent ambient information and a varying
input address different ambiguities.

Another restricted history family sets every chain state initially to
$\gamma u_0$, where $0\leq\gamma\leq1$. For fixed parameters, its response is
exactly $(1-\gamma)T_{\mathrm{cold}}+\gamma T_{\mathrm{equilibrated}}$.
The same interpolation describes a constant pre-window input between zero and
$u_0$ in the pure-delay model. This family does not represent every possible
heater history. Its derivative with respect to $\gamma$ is the difference of
the two endpoint trajectories. If these nuisance sensitivities resemble the
delay sensitivity, fixing history can make delay appear more precise than a
fit that allows history to vary.

More generally, let $A=W^{1/2}J_p$ and $B=W^{1/2}J_a$ be the mean sensitivities
for parameters of interest and nuisance parameters, respectively, under a fixed
positive-definite observation weight $W$. Profiling unrestricted local nuisance
directions leaves information $A^\top(I-BB^+)A$, where $B^+$ is the
Moore--Penrose pseudoinverse. This projection cannot increase information.
Bounds or independent nuisance measurements alter the problem and must be stated.

\subsection{Serial dependence alone does not prove least-squares bias}
For fixed input and known initial conditions, suppose observations have the
correct mean $\bm m(\bm p_0)$ and additive error $\bm\epsilon$ with
$E\bm\epsilon=0$ and covariance $\Sigma$. For any candidate parameter vector
$\bm p$,
\[
 E\|\bm m(\bm p)-\bm m(\bm p_0)-\bm\epsilon\|^2
 =\|\bm m(\bm p)-\bm m(\bm p_0)\|^2+\operatorname{tr}\Sigma.
\]
Correlated errors change sampling uncertainty and efficiency without moving the
minimum of this expected output-error objective. This identity does not imply
an unbiased finite-sample nonlinear estimator. It also does not apply unchanged
when noisy observations enter regressors, the initial observation creates a
nonzero conditional error mean, or the thermal mean is misspecified. The
observed TC-Lab residual offsets and held-out errors require those possibilities
to remain open.

\section{Chain convergence at the causal onset}
\label{sec:onset}
For a cold constant step of magnitude $u_0$ and zero initial thermal deviation,
let $D_n$ have an Erlang
distribution with mean $\theta$ and variance $\theta^2/n$. The temperature
rise relative to ambient is
\[
 y_n(t)=K u_0 E\!\left[
  \left(1-e^{-(t-D_n)/\tau}\right)\mathbf 1_{\{D_n\leq t\}}\right].
\]
At $t=\theta$, the pure-delay response is zero while every finite chain has a
positive response. The central limit theorem and a first-order expansion of
the exponential give
\[
 \sqrt n\,y_n(\theta)\longrightarrow
 \frac{K u_0\theta}{\tau\sqrt{2\pi}}.
\]
To justify the expectation limit, the standardized Erlang variables have unit
second moments, so their positive parts are uniformly integrable. The
exponential remainder is bounded by a constant times
$E(\theta-D_n)^2/\tau^2=O(n^{-1})$.
This is an onset error of order $n^{-1/2}$, even when smooth parts of the
response converge faster. A small average trajectory error can therefore hide
a larger difference near the delay's causal onset. The saved onset calculation
checks the coefficient using an independently tested Erlang-CDF expression.
